光学是一门古老而至今仍十分活跃的学科。现存最早的人工透镜可追溯到公元前七世纪；从以 Newton 为代表的近代物理学开始，光学理论和光学仪器的制造一同进入快速发展时期。今天，光学不仅是基础物理的重要组成部分，也深刻影响着通信、精密测量、医学、天文、材料和信息技术。

贯穿光学发展史的核心问题是：光究竟是什么？牛顿（Isaac Newton）曾把光理解为高速运动的粒子流；同一时期，惠更斯（Christiaan Huygens）主张光是一种波。十九世纪初，杨（Thomas Young）的干涉实验和菲涅尔（Augustin-Jean Fresnel）的理论使波动说重新兴起；麦克斯韦（James Clerk Maxwell）随后指出光是横向电磁波，不同颜色对应不同频率或波长，而真空中的传播速度相同。二十世纪初，爱因斯坦（Albert Einstein）的光量子理论揭示了能量交换的离散性。波粒二象性并不表示单个光子会在“传播时是经典波、探测时是经典粒子”之间切换：光应由量子态或量子化电磁场统一描述，概率振幅可以叠加并产生干涉，探测器则记录一次次离散的能量交换事件。经典波和经典粒子只是不同实验条件下的近似图像；量子电动力学中，光子是电磁场的量子化激发。


光学仪器的发展构成了另一条同样重要的主线。平面镜的平整度、透镜和反射镜的面形精度、材料纯度以及装调误差都会改变仪器的性能。现实中不存在绝对完美的光学系统，设计者必须在功能、精度、口径、重量和成本之间权衡。历史人物与代表性器件提示我们：光的理论与实验仪器始终相互推动。

\historyfigure{optics-history-built-in-v1.png}{光学史人物与代表性实验器件：粒子说、波动说、棱镜色散、望远镜与双缝实验。}

当光学元件和传播尺度远大于光的波长时，波动光学可以取极短波长极限，得到更简单的\textbf{几何光学}。几何光学用“光线”表示能量传播方向，但光线不是有实体的细线，也不是经典粒子的真实轨迹；它是波面法线在空间中的积分曲线。

因此几何光学有明确的适用边界：器件尺寸、曲率半径和传播中被分辨的细节都应远大于波长。接近小孔、尖锐边缘或亚波长结构时，单靠光线会漏掉衍射；两路光重新相遇时，单靠路径也不能判断明暗。后两讲将处理这些波动效应，本讲则先在短波长极限内建立可用于追踪方向的规律。

在任一时刻，把光场中相位相同的点连接起来，就得到等相位面（波面）。光能沿波面的法线方向传播。点光源在均匀各向同性介质中的波面是一组同心球面，所以光线是从光源出发的直线。这给出几何光学的第一条基本规律：\textbf{光在均匀介质中沿直线传播}。

若介质不均匀，波速随位置改变，波面各部分前进速度不同，光线便会弯曲。图 19.1 中，大气密度通常向地面方向增大，折射率也相应增大；接近地平线的太阳光逐层向法线方向偏折。观察者沿最后一段光线反向延长，会把太阳判断在真实位置之上，因此太阳已经略低于几何地平线时仍可能被看见。

\AtmosphereWavefrontPanel

\subsection*{反射}

光照到平整界面时，一部分能量返回原介质，称为反射。在入射点作界面的法线；入射光线与法线的夹角称为入射角 $i$，反射光线与法线的夹角称为反射角 $r$。入射光线、法线和反射光线共面，并且
\begin{equation}
  i=r.
\end{equation}
角度必须相对\textbf{法线}测量。对曲面，应先在入射点作切平面，再以切平面的垂线作为法线。

两块平面镜夹角为 $\alpha$ 时，连续两次反射后，出射方向相对入射方向的总转角为
\begin{equation}
  \Delta\varphi=2\alpha,
\end{equation}
其大小与第一束光的入射角无关。$\alpha=0$ 对应平行镜，可用于平移光路；$\alpha=\pi/2$ 时出射方向与入射方向相反，构成直角反射器。月球激光测距所用的角锥反射器正是这一思想的三维推广。

这个结论可以直接用方向角核验。设一条光线相对固定轴的有向角为 $\varphi$，镜面方向角为 $\beta$。镜面反射把光线关于镜面对称，故一次反射后
\[
 \varphi'=2\beta-\varphi.
\]
若两镜方向角依次为 $\beta_1,\beta_2$，两次反射后
\[
 \varphi''=2\beta_2-(2\beta_1-\varphi)
 =\varphi+2(\beta_2-\beta_1).
\]
因此方向转角确为 $2\alpha$，而不是 $\pi-2\alpha$。后一表达式有时表示两条无向直线间取补角后的夹角，不能当作光的有向传播方向变化。图 19.4 中入射箭头、两次反射点和出射箭头正是按这一关系绘制的。

\historyfigure[0.96\linewidth]{reflection-applications-built-in-v1.png}{图 19.2 与图 19.5：反射定律实验；潜望镜与月面激光反射器。}

\ReflectionPanel

\subsection*{折射}

光从一种透明介质进入另一种透明介质时，传播方向一般会改变。折射光线仍位于入射平面内。若界面两侧折射率分别为 $n_1,n_2$，入射角和折射角分别为 $\theta_1,\theta_2$，则斯涅尔定律（Snell's law）为
\begin{equation}
  n_1\sin\theta_1=n_2\sin\theta_2.
\end{equation}
折射率定义为 $n=c/v$，其中 $c$ 为真空光速，$v$ 为介质中的相速度。在本讲讨论的普通透明介质中 $n>1$。同一介质的折射率通常随波长变化，所以复色光会发生色散。近轴条件 $|\theta|\ll 1$（角度用弧度表示）下，$\sin\theta\simeq\theta$，Snell 定律近似为
\begin{equation}
  n_1\theta_1\simeq n_2\theta_2.
\end{equation}

对顶角为 $A$、置于空气中的三棱镜，记两面外侧入射/出射角为 $i_1,i_2$，棱镜内的折射角为 $r_1,r_2$。由几何关系
\begin{equation}
 r_1+r_2=A,\qquad \delta=i_1+i_2-A.
\end{equation}
在对称光路 $i_1=i_2$、$r_1=r_2=A/2$ 处偏向角取极小值。普通玻璃中蓝光折射率较大，因而偏折通常比红光更强。Newton 的棱镜实验正是由此说明白光由多种颜色组成。

棱镜计算必须把两个界面依次写全。空气折射率取 $1$ 时，第一面与第二面分别满足
\[
 \sin i_1=n\sin r_1,\qquad n\sin r_2=\sin i_2.
\]
给定 $A,n$ 和 $i_1$ 后，先由第一式求 $r_1$，再用 $r_2=A-r_1$ 求 $i_2$，最后才代入 $\delta=i_1+i_2-A$。图示通常取两内部角均为正；若光线越过法线，应使用有向角。令 $r_c=\arcsin(1/n)$，从指定两面贯穿的必要角度条件是 $|r_1|,|r_2|\le r_c$，同时要确认中间光段没有先碰到底面。若第二面上 $|r_2|>r_c$，便发生全反射。等号表示掠射极限，不能把它当作具有有限法向能流的出射光。

把入、出方向交换，光路可逆，因此交换 $r_1,r_2$ 不改变偏向角；但这种对称性本身并不能证明极小。令 $f(r)=\arcsin(n\sin r)$，则 $\delta=f(r_1)+f(A-r_1)-A$。函数 $f'(r)=n\cos r/\sqrt{1-n^2\sin^2r}$ 为偶函数，并随 $|r|$ 增大而增大。由于 $A>0$，$r_1<A/2$ 时有 $|r_1|<|A-r_1|$，故 $\delta'<0$；$r_1>A/2$ 时相反。于是当 $A<2r_c$、对称光路可贯穿棱镜时，唯一极小值在 $r_1=r_2=A/2$，即
\[
 \delta_{\min}=2\arcsin\!\left(n\sin\frac A2\right)-A.
\]
实验中缓慢转动棱镜时，光斑先向一个方向移动、到达最小偏向位置后再反向；这比仅凭一张静态光路图更容易辨认极小值。

\RefractionPrismPanel

图\ref{ch27:fig:deviation}给出固定折射率下的偏向角曲线。入射角增大时偏向角先减小后增大，最低点对应内部对称光路；它与色散图回答的是两个不同问题。
\PrismDeviationDiagram

\subsection*{全反射与光纤}

图 19.8 的 A、B、C 依次为麦克斯韦、爱因斯坦和费曼：他们分别将光与电磁波、光量子和量子电动力学联系起来。D 为全反射实验情景，光从半圆形介质的圆弧面沿半径入射，再在平面边界反射；精确的角度条件以下面的临界角公式判定。
\historyfigure[0.88\linewidth]{optics-history-tir.png}{图 19.8：电磁理论与量子光学的代表人物，以及半圆形玻璃的全反射实验。}

当光由高折射率介质 $n_1$ 射向低折射率介质 $n_2$，折射角随入射角增大得更快。折射角达到 $90^\circ$ 时的入射角称为临界角，满足
\begin{equation}
  \sin\theta_c=\frac{n_2}{n_1},\qquad n_1>n_2.
\end{equation}
当 $i>\theta_c$ 时不再存在传播的折射光，理想界面上发生全反射。

光纤利用芯层 $n_1$ 与包层 $n_2<n_1$ 的界面使光连续全反射，从而让传播方向跟随纤维弯曲。原理虽然简单，但长距离通信要求极低损耗。高锟指出早期光纤的主要损耗来自材料杂质，而非玻璃的不可避免吸收；高纯度材料由此把光纤变成实用通信媒介，他于 2009 年获得 Nobel 物理学奖。

全反射只限制芯层内光线在界面处的角度，并不意味着任意入射光都能被光纤接收。若外界折射率为 $n_0$，子午光线相对光纤轴的最大接收角 $\theta_a$ 满足
\[
 n_0\sin\theta_a=\sqrt{n_1^2-n_2^2}.
\]
右端称为数值孔径。弯曲过急会使外弯侧的界面入射角降到临界角以下，光便从包层泄漏；杂质吸收、散射和弯曲损耗则决定能够传输的实际距离。图 19.9 展示光纤应用，图 19.10 左侧则把芯层、包层和反射位置分开，说明“导光”不是光在单一玻璃棒中无条件反弹。

\historyfigure[0.92\linewidth]{fiber-pioneer-built-in-v1.png}{图 19.9：高锟、低损耗光纤与长距离光通信。}

\FiberFermatPanel

\subsection*{费马原理（Fermat principle）与光程}

逐个界面应用直线传播、反射和折射定律，原则上可以追踪任意光线；但复杂曲面会使法线和面形设计变得困难。Fermat 原理提供了统一表述：真实光路使传播时间相对于邻近光路取驻值。对本讲中的常规问题，它通常表现为极小值。

介质中一段几何路程 $\dd s$ 的传播时间为 $\dd t=n\dd s/c$，因此定义光程
\begin{equation}
  \mathcal L=\int n\,\dd s,\qquad t=\frac{\mathcal L}{c}.
\end{equation}
真实光路使 $\delta\mathcal L=0$。这一表述可推导反射定律和折射定律，也解释为什么光线宁可在低折射率介质中走较长的几何距离。若同一点发出的不同光线经光学系统后严格会聚于同一点，则各条光路的总光程必须相同，否则会聚点处相位不能一致。

以平界面为例，物点与像点到界面的垂直距离分别为 $a,b$，光线在界面上的交点横坐标为 $x$，两点的水平间隔为 $L$。光程
\[
 \mathcal L(x)=n_1\sqrt{a^2+x^2}
 +n_2\sqrt{b^2+(L-x)^2}.
\]
令 $\dd\mathcal L/\dd x=0$，得到
\[
 n_1\frac{x}{\sqrt{a^2+x^2}}
 =n_2\frac{L-x}{\sqrt{b^2+(L-x)^2}},
\]
两边的几何比值分别就是 $\sin i$ 与 $\sin r$，于是恢复 Snell 定律。若两段都在同一介质而界面改为镜面，用镜像点把折线路径展开成直线，便得到 $i=r$。因此直线传播、反射、折射并不是彼此孤立的三条经验规则，而是同一个驻光程原理在不同约束下的结果。

例如，物点 $O=(-s_o,0)$、像点 $I=(s_i,0)$，界面顶点为原点，左侧折射率为 $1$、右侧为 $n$。界面点 $P=(x,y)$ 必须满足
\begin{equation}
 \sqrt{(x+s_o)^2+y^2}+n\sqrt{(x-s_i)^2+y^2}=s_o+n s_i.
\end{equation}
这是一类高次非球面。它说明即使“点到点严格成像”这一基本要求，也往往需要复杂且高精度的曲面；人的角膜与晶状体系统则是自然形成的精密光学系统。

\SphericalMirrorRotatingDiagram

\section{光路计算与费马原理的应用}

\begin{origexamplebox}{例题 19.1【球面镜】}
有一个半径为 $R$ 的球形反射面，圆心位于 $C$，$CV$ 为其对称轴。一束与对称轴平行、距轴为 $d$ 的入射光经球面反射后与对称轴交于 $M$。求 $M$ 与球面顶点 $V$ 之间的距离 $s$。
\end{origexamplebox}
\begin{xesolution}
取子午截面。设入射点为 $P$，令 $u=\sqrt{R^2-d^2}$。从顶点 $V$ 向球心方向量，$P$ 的轴向坐标为 $R-u$。半径 $CP$ 是法线，记它与轴的夹角为 $\alpha$，则 $\sin\alpha=d/R$、$\cos\alpha=u/R$。入射光平行于轴，故反射光与轴的夹角为 $2\alpha$。由图 19.11 左图中的直角三角形，
\[
 s=(R-u)+\frac{d}{\tan 2\alpha}.
\]
又 $\tan2\alpha=2du/(u^2-d^2)$，整理得
\[
 \boxed{s=R-\frac{R^2}{2\sqrt{R^2-d^2}}}.
\]
上式的 $s$ 是从 $V$ 向 $C$ 为正的有向坐标；题目所问几何距离为 $|s|$，参数域为 $|d|<R$。当 $|d|<\sqrt3R/2$ 时 $s>0$，超过此高度时轴交点移到顶点另一侧，$s<0$。当 $d\to0$ 时 $s\to R/2$；不同高度给出不同交点，说明球面镜存在球差。$d=0$ 的轴上光与轴重合，焦点位置须由近轴极限定义。
\end{xesolution}

\begin{origexamplebox}{例题 19.2【转动的镜子】}
一竖直墙上 $A$ 点发出水平光束，平面镜 $OB$ 绕过 $O$ 的水平轴以角速度 $\omega$ 转动。光束经镜面反射后回到墙上形成光点 $D$。已知此刻镜面与墙的夹角为 $\theta$，$AO=d$，忽略光的传播时间，求光点在墙上的移动速度。
\end{origexamplebox}
\begin{xesolution}
在垂直转轴的截面内令墙为 $x=0$，$O=(0,0)$，$A=(0,d)$。镜面直线可写为 $x=-y\tan\theta$，水平入射线 $y=d$ 与镜面交于 $C=(-d\tan\theta,d)$。镜面转过 $\theta$ 时反射方向转过 $2\theta$，故从 $C$ 到墙的竖直位移为 $d\tan\theta\tan2\theta$。于是
\[
 y_D=d+d\tan\theta\tan2\theta.
\]
利用 $\dot\theta=\omega$ 求导，得到沿墙的瞬时速度
\[
 \boxed{v_D=d\omega\!\left(\sec^2\theta\tan2\theta+2\tan\theta\sec^22\theta\right)}.
\]
式子的正负表示移动方向；题目若只问速率，应取绝对值。$\theta$ 接近使反射光平行于墙的位置时，$\tan2\theta$ 发散，光点速度可很大，这并不表示物质或信息以该速度传播。
\end{xesolution}

\CylinderDiagram

\begin{origexamplebox}{例题 19.3【圆柱折射】}
一个半径 $R=2\,\mathrm m$ 的透明圆柱体，其一侧是反射面。一束光从左侧空气入射，经折射和反射后，出射光与入射光平行；两平行光线的距离 $d=2\,\mathrm m$。求透明介质的折射率。
\end{origexamplebox}
\begin{xesolution}
在垂直圆柱轴的截面内分析。图 19.12 中两外部光线关于过圆心 $C$ 的水平轴对称，入射点 $A$ 的高度为 $d/2=R/2$。取 $C=(0,0)$，便有 $A=(-\sqrt3R/2,R/2)$、反射点 $P=(R,0)$，因此相对向内法线 $AC$ 的入射角 $i=30^\circ$。内部线段 $AP$ 相对水平轴向下的倾角为 $\beta$，其中
\[
 \tan\beta=\frac{R/2}{R+\sqrt3R/2}=2-\sqrt3=\tan15^\circ.
\]
所以 $\beta=15^\circ$，折射角为法线与内部光线的夹角 $r=i-\beta=15^\circ$。在 $P$ 处按水平法线反射后，光线到达对称点 $B$，再折射为水平向左的出射光，方向也与题设一致。Snell 定律给出
\[
 1\cdot\sin30^\circ=n\sin15^\circ,\qquad
 \boxed{n=\frac{\sin30^\circ}{\sin15^\circ}=\frac{\sqrt6+\sqrt2}{2}\approx1.93}.
\]
图 19.12 同时画出了两次折射点和右侧反射点；若只凭“平行”而忽略圆柱半径法线的位置关系，角度会误判。
\end{xesolution}

\ShellAndGradientDiagram

\begin{origexamplebox}{例题 19.4【大小球】}
用折射率为 $n$ 的透明物质制成内、外半径分别为 $a,b$ 的空心球，内表面涂有能完全吸收光的物质。平行光射向此球时，球吸收的入射光束横截面积多大？
\end{origexamplebox}
\begin{xesolution}
先采用几何光学模型，忽略外表面的菲涅尔反射损耗，求能到达吸收层的入射光束面积，而不是包含反射损耗的功率吸收截面。取图 19.13 中恰与内球在 $T$ 相切的临界光线。设它进入外球面后的折射角为 $r$，则直线到球心的距离为 $b\sin r=a$。设这条光线在入射平面内的冲量参数（入射线到球心的距离）为 $\rho$，空气侧有 $\rho=b\sin i$。由 Snell 定律 $\sin i=n\sin r$，所以
\[
 \rho=b\sin i=nb\sin r=na.
\]
冲量参数小于 $\rho$ 的光线会碰到吸收层，故
\[
 \boxed{S_{\rm abs}=\pi\min\{b^2,n^2a^2\}}.
\]
若 $na\ge b$，所有射入外球几何截面的光都能到达内层，面积饱和为 $\pi b^2$；这一截断条件不能省略。
\end{xesolution}

\begin{origexamplebox}{例题 19.5【变折射率折射】}
介质位于 $x>0,y>0$，折射率仅随 $y$ 连续变化。一束细光束沿 $x$ 方向垂直入射，并沿半径为 $R$ 的圆弧传播。求 $n(y)$。已知 $n(0)=n_0=1$，材料最大折射率为 $2.5$，求圆弧允许的最大圆心角。
\end{origexamplebox}
\begin{xesolution}
把介质分成许多水平薄层。Snell 定律在各层界面上连续应用，得到守恒量 $n(y)\cos\theta=n_0$，其中 $\theta$ 是光线切线与 $x$ 轴的夹角。对图 19.14 的圆弧，圆心在入射点正上方，有
\[
 \cos\theta=1-\frac{y}{R}.
\]
因此
\[
 \boxed{n(y)=\frac{n_0}{1-y/R}}.
\]
材料条件 $n\le2.5$ 给出 $\cos\theta\ge1/2.5=0.4$，故
\[
 \boxed{\theta_{\max}=\arccos0.4\approx66.4^\circ}.
\]
该角正是从入射半径转到当前半径的圆心角；若把它错当成相对法线的角，会得到互余的错误答案。
\end{xesolution}

\FisherPlateDiagram

\begin{origexamplebox}{例题 19.6【运动学？】}
湖湾由从顶点 $D$ 出发的两条直岸围成，夹角为 $\alpha$，$0<\alpha<\pi$。渔人的房子 $A$ 在一岸的陆地一侧，$C$ 是 $A$ 到该岸的垂足，$AC=h$，$AD=l$，且 $0<h<l$、$C$ 在岸线射线上。朋友的房子 $B$ 关于湖湾角平分线与 $A$ 对称。陆上步行速度为 $v$，划船速度为 $v/2$，可在岸上任意位置上下船。规划从 $A$ 到 $B$ 的最省时路线。
\end{origexamplebox}
\begin{xesolution}
先分清陆地和水域。图 19.15 中 $DC$ 是一侧岸线，$A,B$ 都不在水中；因此不能从房子直接划船。令
\[
 c=DC=\sqrt{l^2-h^2},\qquad \beta=\alpha/2.
\]
两岸上船、下船点记为 $P,Q$，$DP=t_1,DQ=t_2$。直线步行段均在各自陆地半平面内，船的直线段在湖湾内。其时间为
\[
 vT=\sqrt{h^2+(t_1-c)^2}+\sqrt{h^2+(t_2-c)^2}
 +2\sqrt{t_1^2+t_2^2-2t_1t_2\cos\alpha}.
\]
为什么可以取对称路线？前两项是严格凸函数，第三项是线性坐标差的欧氏范数，也是凸函数；交换 $t_1,t_2$ 不改变时间。将一条路线与它的对称路线取参数平均，凸性给出
$T((t_1+t_2)/2,(t_1+t_2)/2)\le T(t_1,t_2)$。
所以最优点必有 $t_1=t_2=t\ge0$。这不是仅凭图形对称猜测。
此时 $PQ=2t\sin\beta$，有
\[
 T(t)=\frac2v\left[\sqrt{h^2+(t-c)^2}+2t\sin\beta\right],
 \quad
 \frac v2T'(t)=\frac{t-c}{\sqrt{h^2+(t-c)^2}}+2\sin\beta.
\]
令 $k=2\sin\beta$，且
$T''(t)=2h^2/[v(h^2+(t-c)^2)^{3/2}]>0$。
由于 $T'(0)=2(k-c/l)/v$，答案完整分为
\[
 \boxed{\begin{cases}
 k\ge c/l:&t_*=0,\quad T_{\min}=2l/v;\\[2pt]
 0<k<c/l:&t_*=c-\dfrac{kh}{\sqrt{1-k^2}},\quad
 T_{\min}=\dfrac2v\left[h\sqrt{1-k^2}+kc\right].
 \end{cases}}
\]
第一种情形沿陆地 $A\to D\to B$ 行走，不划船；第二种情形沿 $A\to P$ 步行、$P\to Q$ 划船、$Q\to B$ 步行。
后一条件自动保证 $k<1$ 和 $t_*>0$。两式在 $k=c/l$ 处连续，时间单位均为长度除以速度；$h$ 不由 $l\sin\alpha$ 决定。
同岸的多次上下船不会更快，因为把这段船路改为沿同岸直线步行既不更长、速度又更大。于是上述一次横渡或完全绕岸的比较已覆盖最优路线。
它体现了费马最短时间思想：总距离最短不等于总时间最短。
\end{xesolution}

\begin{origexamplebox}{例题 19.7【严格聚焦】}
球面镜不能把所有平行于轴的入射光严格会聚于一点。能够完成这一功能的反射面应具有何种形状？
\end{origexamplebox}
\begin{xesolution}
取顶点为原点、轴向右为 $x$ 正方向，焦点 $F=(f,0)$，$f>0$。平行光必须从右向左入射，照到向右开口的反射面。参考平面取 $x=a$，位于整个镜面右侧；反射点为 $P=(x,y)$。轴上光路长为 $a+f$，等光程条件为
\[
 (a-x)+\sqrt{(x-f)^2+y^2}=a+f.
\]
故 $PF=x+f$。平方并消去 $x^2+f^2$，得到 $y^2=4fx$；它也表明 $PF$ 等于 $P$ 到准线 $x=-f$ 的距离。因此子午截面是抛物线，旋转后为抛物面：
\[
 \boxed{y^2=4fx}
\]
也可直接检验反射定律。曲面法向可取 $\boldsymbol N=(-2f,y)$，入射单位方向为 $\boldsymbol k=(-1,0)$。反射方向
\[
 \boldsymbol k'=\boldsymbol k-2\frac{\boldsymbol k\cdot\boldsymbol N}{\boldsymbol N\cdot\boldsymbol N}\boldsymbol N
 =\frac{(4f^2-y^2,-4fy)}{4f^2+y^2}
 =\frac{F-P}{|F-P|}.
\]
所以反射光真正射向 $F$，而非仅延长线经过它。抛物面反射镜对轴上平行光严格无球差；有限口径带来衍射，制造误差和离轴入射仍会带来像差。
\end{xesolution}

\begin{origexamplebox}{例题 19.8【圆形光线】}
一颗半径为 $R$ 的行星，其大气折射率随高度 $h$ 线性变化：$n=n_0-\alpha h$，其中 $n_0$ 是表面折射率。若某高度上沿行星表面切向发出的光恰能绕行星作圆周传播，求该高度 $h_0$。
\end{origexamplebox}
\begin{xesolution}
球对称介质的射线守恒量为 $nr\sin\psi$，$\psi$ 是光线与径向的夹角。圆形轨道要求 $\psi=90^\circ$，并且 $n(r)r$ 在该半径处取驻值：
\[
 \frac{\dd}{\dd r}[n(r)r]=0,\qquad r=R+h.
\]
代入 $n=n_0-\alpha h$ 得
\[
 n_0-\alpha h-\alpha(R+h)=0,\qquad
 \boxed{h_0=\frac{n_0-\alpha R}{2\alpha}}.
\]
结果还必须满足 $h_0\ge0$ 且 $n(h_0)>0$；若材料参数不满足这些条件，物理解不存在。
\end{xesolution}

\begin{origexamplebox}{例题 19.9【平板玻璃砖】}
观察者距折射率为 $n$、厚度为 $D$ 的平行玻璃板近表面为 $L$。发光点位于玻璃板另一侧，距远表面为 $S$。求观察者看到的发散光线反向延长后的视距离。
\end{origexamplebox}
\begin{xesolution}
取一条近轴光线。在空气中角度为 $\theta$，玻璃中角度为 $\theta_g$。由近轴 Snell 定律 $\theta_g\simeq\theta/n$。光线在三个区段产生的横向位移为
\[
 y\simeq L\theta+D\theta_g+S\theta
  =\theta\left(L+\frac Dn+S\right).
\]
观察者按空气中的末段方向反向延长，虚像点到观察者的距离就是 $y/\theta$：
\[
 \boxed{L_{\rm view}=L+\frac Dn+S}.
\]
所以玻璃的视厚度为 $D/n$，比真实厚度小。该结果依赖近轴条件；大角度光线的反向延长线不再严格交于同一点。
\end{xesolution}

\begin{origexamplebox}{例题 19.10【全反射】}
为什么从玻璃砖的一个侧面通常看不到由相邻侧面射入砖内的光线？玻璃折射率最大为多少时，才可能从侧面看见相邻侧面入射的光？
\end{origexamplebox}
\begin{xesolution}
设光从空气经一个面进入玻璃，玻璃内相对该面法线的折射角为 $r$。相邻面法线与原法线垂直，所以光到相邻面时入射角为 $90^\circ-r$。要从相邻面射出，必须满足
\[
 90^\circ-r\le\theta_c,\qquad \sin\theta_c=\frac1n.
\]
而空气侧入射角最大为 $90^\circ$，故玻璃内最大折射角满足 $\sin r_{\max}=1/n$，即 $r_{\max}=\theta_c$。存在可出射光线的条件是 $90^\circ-\theta_c\le\theta_c$，从而
\[
 \theta_c\ge45^\circ,\qquad \boxed{n\le\sqrt2}.
\]
这里的等号只是双掠射极限：$n=\sqrt2$ 要求入射和出射的空气侧光线都沿表面掠行，不能提供有限的贯穿通光量。因此实际非零通光量的条件为 $1<n<\sqrt2$，上确界为 $\sqrt2$。常见玻璃 $n\approx1.5>\sqrt2$，相邻面上发生全反射，不能沿该路径直接看到入射光束。
\end{xesolution}

\begin{origexamplebox}{例题 19.11【最小偏向角】}
一三棱镜顶角为 $A$、折射率为 $n$，光由一个面入射并从另一面射出。求偏向角的可能范围。
\end{origexamplebox}
\begin{xesolution}
取空气中 $n>1$、$0<A<\pi$，讨论从指定两侧面贯穿且未被底面截断的光路。为覆盖小顶角情形，内部角 $r_1,r_2$ 采用有向角，满足 $r_1+r_2=A$；负值表示光线越过对应法线。外侧角为 $i_j=\arcsin(n\sin r_j)$，偏向角 $\delta=i_1+i_2-A$。
令 $r_c=\arcsin(1/n)$，两面透射要求 $|r_j|\le r_c$，于是
\[
 A-r_c\le r_1\le r_c.
\]
若 $A>2r_c$，区间为空，没有贯穿两面的透射光路；若 $A=2r_c$，只剩双掠射极限。以下取 $0<A<2r_c$。
写 $F(r)=\arcsin(n\sin r)$，其导数
\[
 F'(r)=\frac{n\cos r}{\sqrt{1-n^2\sin^2r}}
\]
为偶函数，随 $|r|$ 严格增加（$n>1$）。因此
$\delta'(r_1)=F'(r_1)-F'(A-r_1)$ 在 $r_1<A/2$ 时为负，在 $r_1>A/2$ 时为正。
最小值确实在 $r_1=r_2=A/2$，而不只由可逆性猜出：
\[
 \boxed{\delta_{\min}=2\arcsin\!\left(n\sin\frac A2\right)-A}.
\]
偏向角上界出现在某一外侧光线掠射界面。取 $r_1=r_c$，$r_2=A-r_c$，得到
\[
 \boxed{\delta_{\max}=\frac\pi2+\arcsin[n\sin(A-r_c)]-A}.
\]
这一式在 $A<r_c$ 时仍成立，此时第二个反正弦为负，不能擅自截为零。实际非掠射光路的范围为 $\delta_{\min}\le\delta<\delta_{\max}$；若允许极限光线，则右端也可写等号。特定入射位置、有限底面和光阑还会缩小可达范围，以上是可选择入射位置的两侧面几何范围。角度统一用弧度。
\end{xesolution}

\RainbowDiagram

\begin{origexamplebox}{例题 19.12【彩虹成因】}
雨后出现彩虹。把空中的水滴看作折射率为 $n$ 的球，解释主虹的形成，并求其特征方向。
\end{origexamplebox}
\begin{xesolution}
主虹对应“入射折射—一次内反射—出射折射”的光路，如图\ref{ch27:fig:rainbow}左侧路径。入射点 $A$ 的半径是法线。设空气侧入射角为 $i$、水滴内折射角为 $r$，$\sin i=n\sin r$。光线相对入射方向的总偏折为
\[
 D_1=\pi+2i-4r.
\]
大量邻近光线在 $\dd D_1/\dd i=0$ 的方向聚集，因而亮度最大。由 $\cos i=n\cos r\,\dd r/\dd i$ 得
\[
 2-4\frac{\cos i}{n\cos r}=0
 \quad\Longrightarrow\quad
 \sin^2i=\frac{4-n^2}{3}.
\]
对水的可见光折射率代入，主虹的观察半径约为 $42^\circ$。折射率随波长变化，红光偏折较小，故主虹红色在外、紫色在内。单个水滴只把某一颜色送入某一方向；整条虹来自视线方向不同的大量水滴。
\end{xesolution}

\begin{origexamplebox}{例题 19.13【广义彩虹】}
网页和自然界中还能看到多种类似彩虹的彩色弧带或光环。怎样判断并解释这些图样的成因？
\end{origexamplebox}
\begin{xesolution}
不能把所有彩色圆弧都归结为同一种几何光路，应先看光源方向、弧心位置、颜色次序和介质形态。
\begin{itemize}
  \item 次虹由水滴内两次反射形成，偏折 $D_2=2\pi+2i-6r$；它位于主虹外侧、亮度较弱，颜色次序与主虹相反，主虹与次虹之间形成较暗的 Alexander 带。
  \item 雾虹仍与水滴有关，但液滴很小，衍射展宽压过色散，因而弧带宽而近白色。
  \item 日晕和月晕主要由高空六角冰晶的折射形成，常见半径约为 $22^\circ$ 或 $46^\circ$；其圆心在太阳或月亮，而普通虹的圆心在反太阳点。
  \item “宝光”或 glory 出现在反太阳方向，涉及小液滴的后向散射、表面波和衍射，不能仅用一次或两次几何反射解释。
\end{itemize}
因此，几何光学能够给出许多弧带的主方向和色序，但小尺度粒子造成的展宽、干涉和衍射必须留到波动光学中处理。
\end{xesolution}
